% titel: Blatt 9, A7: Kc, Kp und Le-Chatelier für eine exotherme Reaktion
% skript: Mortimer Kap. 14 (Gleichgewicht, Kp, Le Chatelier)

% Teil I: Werkzeuge, kein Lösungsweg.

An einer konkreten gasförmigen Reaktion soll das Massenwirkungsgesetz in
Konzentrationen und Drücken formuliert und die Verschiebung des
Gleichgewichts vorhergesagt werden.

\begin{werkzeugkasten}[MWG, Kp und Le Chatelier]
Für die Gasreaktion $2H_2S + CH_4 \rightleftharpoons CS_2 + 4H_2$:
\begin{itemize}
  \item $K_c$ aus den Gleichgewichtskonzentrationen mit Koeffizienten als
        Exponenten.
  \item $K_p$ analog mit den Partialdrücken der Gase.
  \item Zusammenhang: $K_p = K_c(RT)^{\Delta n}$, wobei $\Delta n$ die
        Differenz der Gaskoeffizienten (Produkte $-$ Edukte) ist.
  \item Temperaturerhöhung begünstigt die endotherme Richtung; hier ist die
        Reaktion nach rechts \emph{exotherm}, also begünstigt
        Temperaturerhöhung die \emph{linke} Seite.
  \item Druckerhöhung begünstigt die Seite mit weniger Gaspartikeln;
        betrachte die Teilchenzahlen links und rechts.
\end{itemize}
Dynamik: Zugabe eines Edukts (H$_2$S) verschiebt nach rechts; ein
Katalysator ändert die Lage nicht.
\textit{Skript:} Mortimer Kap. 14.
\end{werkzeugkasten}

\begin{achtung}[Darauf achten]
\begin{itemize}
  \item Gaskoeffizienten zählen: links 2$+$1$=$3 Teilchen, rechts 1$+$4$=$5.
  \item Da rechts mehr Gasmoleküle sind, verschiebt Druckerhöhung nach
        \emph{links}.
  \item $K_p = K_c(RT)^{\Delta n}$ mit $\Delta n = (1+4)-(2+1) = 2$.
  \item ``Exotherm nach rechts`` $\Rightarrow$ bei Temperaturerhöhung
        Gleichgewicht nach links (kleinstes Zwangprinzip).
\end{itemize}
\end{achtung}
